\scanpage{344}
\mspage{44}

celle-ci est le plus souvent non moins cachée dans l’inconscient que celle-là, tout en s’exprimant de façon tout aussi flagrante dans vos actes et dans vos comportements.

\noindent 19. La violence de rétaliation et les préceptes de Jésus

Il me reste enfin à parler encore d’un type d’agressivité, ou de violence, qui n’est qu’indirectement liée aux pulsions d’antagonisme vis-à-vis des archétypes fondamentaux. Il s’agit de ce qu’on pourrait appeler la ``\textit{violence de réaction}''. C’est la violence qui naît dans la psyché en réaction à des agressions subies, réaction d’autant plus forte que ces agressions sont ressenties comme injustifiées (par exemple, comme excessives), voire même comme totalement ``gratuites''.

La violence de réaction prend souvent la forme d’une \textit{rétaliation}, s’exerçant à l’encontre de l’agresseur (réel ou présumé). Encore faut-il que l’agresseur soit connu et qu’il soit vulnérable. Il n’en est rien pour les ``\textit{agressions diffuses}'', c’est-à-dire pour celles qui ne sont pas le fait de telle personne en particulier, mais de tout un milieu ou de conditions de vie, reconnues par les couches profondes comme incompatibles avec le bien-être et l’épanouissement de la psyché. Les agressions subies tout au cours de l’enfance de la part du milieu parental sont le plus souvent de cet ordre. Un autre cas fréquent où une libération des tensions psychiques par une violence de retaliation est freinée, et souvent même réduite à des actes purement symboliques et de faible effet sur la psyché, est celui où l’agresseur est ressenti (à tort ou à raison) comme trop puissant pour qu’une agression contre lui soit pensable. (Si ce n’est tout au plus dans des occasions propices, quand il est ressenti comme étant en position de faiblesse, et quand l’âme croit pouvoir frapper sans risques...)

Mais alors même que l’agressivité de réaction accumulée par la psyché ne pourrait trouver exutoire et se libérer par des actes de rétaliation, elle trouve nécessairement une expression erratique, par ce qu’on peut appeler la ``\textit{violence de défoulement}''. C’est encore une forme de violence gratuite, frappant au gré des occasions partout où elle peut frapper avec un minimum de risques de rétaliation. Cette forme de violence prolifère surtout en périodes de crise, ou dans des sociétés en voie de décomposition. Elle fait partie des nombreux ``signes de la Fin'', qui marquent de plus en plus au fil des ans et des décennies, le ``\textit{Temps de la Fin}'' dans lequel nous sommes engagés depuis la première guerre mondiale, et qui va toucher à sa fin le 14 Octobre 1996.

Les mécanismes de la ``violence de réaction'', par quoi toute violence subie suscite une violence en retour, fait partie des \textit{conditionnements originels de l’âme}. La libération de la dépendance par rapport aux Figures archétypes, laquelle met fin à la ``violence automatique'' (ou de transfert) liée à ces figures, ne met pas fin pour autant à ce conditionnement plus tenace et plus impérieux encore que cette dépendance. Comme pour tout autre conditionnement originel, la ``vulnérabilité à l’agression'' ne prend fin que par un Acte de Dieu, entérinant les efforts faits par l’âme pour \unsure{le} dépasser. Le plus souvent, cet Acte de Dieu s’accomplit par l’Acte de libération. (Cet Acte met fin à chacun des conditionnements originels qui restent actifs dans l’âme, à la seule exception de la répression sexuelle.)

De cette vulnérabilité de l’âme à la violence, la plupart des messagers du Nouvel Age (et parmi eux la plupart des destinataires de la Lettre) ne seront pas libérés dès leur présente existence, mais seulement au terme

% archive sheet number: 344
